\documentclass{article}
\usepackage{algorithm, algorithmic, amsfonts, amsmath}

\begin{document}
\begin{algorithm}[!h]
\renewcommand{\algorithmicrequire}{\textbf{Input:}}
\renewcommand{\algorithmicensure}{\textbf{Output:}}
	\caption{PTAN (positive tree augmented naive Bayes), with the averaged variant APTAN}
	\begin{algorithmic}[1]
	    \REQUIRE
	    Positive set $E_{1}$ ($N_{E_{1}}$ cases); unlabeled set $E_{u}$ ($N_{E_{u}}$ cases);
	    predictors $\mathbf{X}=\{X_{1},\dots,X_{n}\}$ with $X_{i}$ taking $r_{i}$ values;
	    class prior $p=P(C=1)$
	    \ENSURE
	    Tree structure and all conditional probability parameters
	    \STATE \textbf{Key identity:} the negative-class parameters are \emph{not} estimated from
	    data --- they are recovered from the unlabeled marginal by subtracting the positive
	    contribution weighted by $p$ % this step is the defining operation of branch C: an assumed-known prior placed inside a probabilistic model
	    \STATE \textbf{Stage 1: structure learning}
	    \STATE Estimate $P(x_{ij}\mid 1)$ from $E_{1}$ with Laplace smoothing
	    \STATE $R_{i}(j)\leftarrow N_{ij}^{E_{u}}-P(x_{ij}\mid 1)\,p\,N_{E_{u}}$;\quad
	    normalise by $Z_{i}$ and smooth to obtain $P(x_{ij}\mid 0)$ % the negative marginal is U minus the positive contribution, truncated at zero and renormalised
	    \STATE Repeat the same construction on \emph{pairs} to obtain
	    $P(x_{ij},x_{kl}\mid 1)$, $P(x_{ij},x_{kl})$ and $P(x_{ij},x_{kl}\mid 0)$
	    \STATE Compute the conditional mutual information $I(X_{i},X_{k}\mid C)$, whose two terms
	    are weighted by $p$ and by $\big(P(x_{ij},x_{kl})-P(x_{ij},x_{kl}\mid 1)p\big)$ respectively
	    \STATE Build the maximum spanning tree over $\mathbf{X}$ with edge weights
	    $I(X_{i},X_{k}\mid C)$; root it and orient the edges to fix each tree parent $Pa^{\ast}(X_{i})$ % Friedman's TAN structure learning, used unchanged
	    \STATE Set $Pa(X_{i})=\{C\}\cup Pa^{\ast}(X_{i})$
	    \STATE \textbf{Stage 2: parameter learning}
	    \STATE Estimate $P(x_{ij}\mid 1,pa_{i}^{\ast})$ from $E_{1}$ with Laplace smoothing
	    \STATE Apply the same subtraction identity conditioned on $pa_{i}^{\ast}$ to obtain
	    $P(x_{ij}\mid 0,pa_{i}^{\ast})$ % the same identity applied again, locally, at each value of the tree parent
	    \STATE \textbf{APTAN variant:} $p$ is rarely known exactly, so treat it as
	    $p\sim\mathrm{Beta}(a,b)$ and integrate it out
	    \STATE \quad Negative-class probabilities take the resulting closed form; negatives are
	    replaced by $1/r_{i}$ (resp.\ $1/(r_{i}r_{k})$) and renormalised % the integrated estimate can be negative, so it is truncated and renormalised
	    \STATE \quad In the mutual-information term, $P(x_{ij},x_{kl}\mid 1)\,p$ is replaced by
	    $P(x_{ij},x_{kl}\mid 1)\cdot\dfrac{a}{a+b}$ % the expectation of a Beta distribution replaces a fixed p
		\RETURN the PTAN (or APTAN) classifier
	\end{algorithmic}
\end{algorithm}
% Calvo, Larrañaga, Lozano, Pattern Recognition Letters 2007,
% "Learning Bayesian Classifiers from Positive and Unlabeled Examples".
% Position: extends positive naive Bayes (PNB) to a tree-augmented naive Bayes (TAN), allowing one layer of dependence
% among the predictors instead of full conditional independence.
% Why branch C: the defining operation of the whole paper is the identity on line 4, in which none of the negative
% probability parameters is estimated from data. Each is the unlabeled marginal minus the positive contribution
% weighted by p, renormalised. The structure learning, conditional mutual information plus a maximum spanning tree,
% follows Friedman's TAN unchanged. APNB and APTAN address the case of unknown p: p is treated as a random variable
% with a conjugate Beta(a,b) prior and integrated out over [0,1], which gives a closed-form estimate.
% The cost is that the integrated estimate can turn negative and must be truncated to a small value and renormalised.
%
% Note: this pseudocode follows the concise version. The smoothing term and the full APTAN closed form are in the
% original, and both are worth checking against it before implementing.
\end{document}
